\section{Method}
\label{sec:method}

The system has four parts: a memory store that records what was observed and when
(Section~\ref{sec:memory}), a world model that turns interaction histories into typed latent estimates and
outcome predictions (Section~\ref{sec:model}), planners that choose skill parameters and decide when to probe
(Section~\ref{sec:planner}), and an adaptation of a pretrained looped language model that hosts the same
interface (Section~\ref{sec:inkling}).

\subsection{Typed, versioned, causal memory store}
\label{sec:memory}
% Source: docs/findings.md "what was built"; ENVS.md "Observations".

The record layer is a property graph over entities (objects, robot links, regions). Each observation is an
immutable \emph{evidence event} carrying an entity, a typed property (pose, wrench, joint, gripper state,
tactile, size, shape, region, and so on), a value with units and frame, a source, and two timestamps: when
it was measured and when it becomes available to the actor. The store provides:
\begin{itemize}
  \item \textbf{versioning}: revisions are explicit; nothing is overwritten in place;
  \item \textbf{causal availability}: late packets are buffered and become readable only at their availability
        time, so a query at time $t$ never sees the future;
  \item \textbf{evidence hygiene}: duplicate suppression, same-source correlation (repeated readings from one
        source do not inflate confidence), authoritative correction, and retraction or deletion that removes
        every readable view;
  \item \textbf{provenance}: estimates written back by a model are marked as model-written and derived;
  \item \textbf{an actor boundary}: evaluator-only labels (true hidden physics, true contact graph) are rejected
        if anything on the actor side tries to read them.
\end{itemize}
Several invariances hold exactly by construction at this layer (duplicate and replay invariance, opaque-identifier
renaming, event-order permutation, time-origin shift, unit conversion, frame round trips, quaternion sign).
Latent physical properties such as grip friction, restitution, compliance, clearance and joint friction are part
of the property vocabulary but are never emitted by a sensor; they exist only as model-written estimates.

\subsection{World model with a typed-slot bottleneck}
\label{sec:model}
% Source: docs/MODEL_V2.md sec. 2, 10a, 11a.

The sprint models (Section~\ref{sec:res-sprint}) let history reach the prediction head through several routes,
and the typed slot was ignored. The second-generation model (Figure~\ref{fig:arch}; 2.28\,M parameters in the
typed variant, the generic variant within 10\%) removes the bypass.

\paragraph{History path.} Up to six past interactions with the target are encoded. Each interaction is a
$64\times26$ resampled time-series segment of actor-legal signals (three strided 1-D convolutions, two
transformer layers, attention pooling), the skill call that produced it, 14 analytic features, and an embodiment
vector. Each interaction yields a diagonal Gaussian over the typed latents (ten in the first version; a crush
limit was added in round~3).

\paragraph{Fusion.} Per-interaction estimates are fused by precision weighting with a learned prior conditioned
on the target and embodiment tokens, with a learned equicorrelation correction so that $n$ correlated readings
count for less than $n$ independent ones:
\begin{equation}
  \mathrm{prec}_\ell \;=\; p_{0,\ell} \;+\; \sum_{k=1}^{n} \frac{\mathrm{prec}_{k,\ell}}{1 + \rho_\ell\,(n-1)} .
\end{equation}
The result is one \emph{typed slot} (mean and variance, in normalised physical units) per latent. Slots can be
read out, and overwritten for interventions. A per-latent variance temperature is fitted on validation data.

\paragraph{Scene path.} The current entity state (1 to 256 entity tokens) goes through an MLP and slot attention
with eight slots \citep{locatello2020slot}, is joined by target and embodiment tokens, and is reduced by two
transformer layers to a context vector of width 128. The scene path sees only the current state.

\paragraph{Head.} A FiLM-conditioned residual MLP \citep{perez2018film} takes the candidate skill call, the
context and the slot values, and predicts eight continuous outcomes with heteroscedastic variance, event logits
(success, slip, drop, jam, collision, stall; damage from round~3), and log effort. Physics margins are given
to the head explicitly, for example the crush margin $\log(\text{crush}) - \log(\text{grip force})$ and a
slip margin from mass and grip friction, so that the minimal safe grip is an explicit trade-off.

\paragraph{Non-bypass property.} Interaction histories reach the head \emph{only} through the typed slots. A
unit test checks that with all slots overwritten, predictions are identical with and without history and change
when the slots change.

\paragraph{Training.} System identification is trained with per-interaction, fused-posterior and prior negative
log-likelihoods against evaluator labels; the outcome loss is detached from the estimator. The head receives,
per sample, either the true latent (with a probability annealed from 0.75 to 0.5) or a sample of the fused
posterior, always presented as a known value, which is what the planner feeds at test time. History dropout is
0.15. Optimisation is AdamW \citep{loshchilov2019adamw} with a one-cycle schedule and family-homogeneous batches.

\paragraph{Variants.} \emph{Typed} (above); \emph{generic} (same budget, an unstructured memory vector instead
of typed slots); \emph{no-memory} (no interaction history); \emph{oracle} (the typed model with the true
latents written into the slots; always labelled, an upper reference rather than a competitor).

\begin{figure}[!htbp]
  \centering
  \fbox{\parbox{0.92\linewidth}{\small\ttfamily\raggedright
  past interactions ($\le 6$): segment, skill call, analytic features, embodiment\\
  \hspace*{1.5em}$\rightarrow$ per-interaction Gaussian over typed latents\\
  \hspace*{1.5em}$\rightarrow$ precision fusion with learned prior $\rightarrow$ TYPED SLOTS (mean, var)\\[2pt]
  current entity state: entity MLP, slot attention, target + embodiment tokens\\
  \hspace*{1.5em}$\rightarrow$ transformer $\times 2$ $\rightarrow$ context $c$\\[2pt]
  slot values $z$, candidate skill call, $c$ $\rightarrow$ FiLM head\\
  \hspace*{1.5em}$\rightarrow$ continuous outcomes, event probabilities, effort}}
  \caption{World model, schematic. History reaches the head only through the typed slots.
  \todo{replace this text box with a drawn architecture figure}}
  \label{fig:arch}
\end{figure}

\subsection{Planners}
\label{sec:planner}
% Source: docs/MODEL_V2.md sec. 3, 9g, 10a, 11a, 11e.

\paragraph{Sampling planner.} The base planner runs the cross-entropy method \citep{rubinstein1999cross}
(population 256, four iterations, 10\% elites) over the scalar parameters of the candidate skill calls proposed
by the environment. The cost of an action under latent sample $s$ combines failure, missing progress, effort
relative to a state-dependent reference, risk and duration; the score is the mean over 32 posterior samples of
the slots plus a multiple of their standard deviation. A value-of-information term (expected value of perfect
information, scaled by the fraction of decision-relevant variance a probe is predicted to remove) decides
whether to spend a probe.

\paragraph{Residual planner with a competence gate.} Planning freely with an imperfect model was the main source
of closed-loop losses (Section~\ref{sec:res-r4r6}). The residual planner therefore starts from the scripted
call and changes only its continuous parameters, never the discrete choice of skill and target, and only when
a lower confidence bound on the improvement, computed from a three-seed ensemble and the latent posterior, is
positive. A \emph{competence gate} allows deviations only in families whose validation success AUC is at least
0.8; elsewhere the planner plays the scripted action. A variant is anchored on the \pta{} script instead of the
nominal one.

\paragraph{Lot planner.} For lots of parts that share hidden physics, interactions with every part are fused
under one lot identifier. For each pick the planner scores grip force $\times$ lift speed candidates by
predicted damage, slip and drop risk plus a time cost. The pre-registered rule (round~4) probes when the upper
bound on the best candidate's risk exceeds a threshold $\tau$ tuned on validation. The later \emph{crush-gated}
rule uses the typed posterior directly: probe only while the lot's crush posterior is still wide and the part
is unprobed, and cap the grip at 0.7 times the posterior 10th percentile when the lot is believed fragile.

\subsection{\sysI{} and \sysII{} on a pretrained looped transformer}
\label{sec:inkling}
% Source: docs/INKLING_SYS12.md sec. b; docs/findings.md "inkling gate".

We adapt the pretrained \inkling{} language model, a looped transformer whose core block is applied three times
\citep{dehghani2019universal}, at a pinned revision with verified weight checksums.

\paragraph{\sysII{} (slow).} \sysII{} is the base language model with every adapter off. It reads an
actor-legal summary of the cell state and the list of allowed calls, and selects one by constrained
log-likelihood over handoff lines of the form \texttt{<sys1> skill=<skill> target=<entity>}. A gate test
checks that this path is bit-identical to the unmodified model.

\paragraph{\sysI{} (fast).} After the \sysII{} context, a trained mode-token embedding, one action token (the
skill chosen by \sysII) and 16 readout tokens are appended. \sysI{} runs the pre-loop layers and \emph{one} of
the three native core passes; post-loop layers and the language head are skipped. Typed memory enters only
through zero-initialised, gated set cross-attention \citep{lee2019set,jaegle2021perceiver}: memory to the
\sysI{} positions at the start of the pass (read), and \sysI{} positions to memory at the end (write).
Low-rank adapters \citep{hu2022lora} of rank 8 on the query, key, value and output projections of the twelve
core layers (48 modules) have three modes: \emph{off} (\sysII), \emph{gated} (active only at positions from the
mode token onward, so every \sysII{} prefix position keeps its base state and its key-value cache can be reused),
and \emph{always} (an ungated control). The readout feeds the same outcome head as Section~\ref{sec:model}, so
the same planners run on top. The memory tokens come from the frozen typed front end of the world model:
target, embodiment, eight entity slots, and eleven typed-latent tokens (posterior mean, log-variance, evidence
count).

\paragraph{Memory footprint.} The model's quantised expert weights are kept packed on the GPU and only the
experts a batch routes to are dequantised on the fly, reproducing the reference arithmetic bit for bit. This
brings the footprint from 52.8\,GB (dequantised) to 20.6\,GB allocated (22.2\,GB peak), and \sysI{} training at
batch size 32 stays near 25\,GB.

\paragraph{Controls.} Ungated adapters; the same architecture with shuffled core weights (matched weight
statistics, destroyed structure); a compact two-layer transformer \sysI{} over the same memory tokens; and the
world model of Section~\ref{sec:model} trained with the same number of samples.
